\section{Manuscript updates and a further research program}
\label{sec:agenda}

The four parts of this article have different mathematical outputs.
The Euler theorem proves a uniform inequality and identifies its exact
optimal constant. The golden certificates prove special-value and
functional identities. The radial theorem classifies a local change
in the shape of a zero curve. The Cayley theorem constructs an
algebraic quotient and separates a target from a specific relation
ideal. Their common feature is that each conclusion has a finite,
explicit point of verification, together with an analytic or algebraic
argument explaining why that verification suffices.

\subsection{Corrections of status and scope}

No false displayed formula was established in the source material
audited for this continuation. The proposed corrections are therefore
specific updates of proof status and precise limitations on how the
new results may be used. A blanket assertion that the entire
375-page manuscript has been verified would exceed this audit.

\paragraph{The universal Euler question is settled.}
The question in \src{chapters/10-discovery.tex} asking for a universal
real-order Euler bound above the critical triangle can be replaced by
\cref{eul:thm:main}. Its earlier parameter-dependent bounds remain valid
and can still be useful when they are smaller. The new result provides
the universal budget $9/(8\,2^N)$ for $N\ge2$ and identifies the optimal
constant for all $N\ge1$. The exact critical-triangle constant also
remains a separate sharper statement on its restricted parameter set.

The source's warning that the Gaussian maximum alone does not furnish
the universal Euler bound is mathematically correct. It should be
retained as an explanation of the missing implication, followed by the
new positive-tail representation and uniform kernel gap. The proof
cannot be shortened by deleting that step. In particular, a bound on
the complete sum does not automatically bound every rescaled tail.

\paragraph{The remaining displayed golden weight-five formulas are proved.}
The numerical-evidence paragraph preceding the second and third
weight-five formulas in \src{chapters/03-algebraic.tex} can now cite
\cref{gold:thm:main}. The canonical index-$12$, $20$, and $24$
normalizations and their lower-weight vanishing statements can be
updated using the exact descent and recombination in
\cref{gold:sec}. The earlier numerical checks are useful historical
evidence but are no longer the justification for these identities.

The weight-$6$ through weight-$9$ formulas retain their prior
conjectural status. Contracting a weight-five tensor proves identities
of lower weight; it cannot create a weight-six tensor identity.
Likewise, a successful finite alphabet does not prove that its
specializations exhaust all numerical relations at the golden ratio.

\paragraph{The quartic-transition question has a local answer.}
The decimal transition in \src{chapters/05-real-turning.tex} can be
supplemented by the exact rational rectangle in
\cref{turn:thm:transition}. The higher radial behavior is given by the
negative sextic coefficient and \cref{turn:thm:fold}. There is exactly
one transition in the stated rectangle, and an open local region with
two small-radius turning points. Neither assertion proves that this is
the only transition on the entire threshold curve or that these are
the only turning points in the full disk.

\paragraph{The Cayley obstruction is strengthened, while $S_6$ stays open.}
The projector in \cref{cay:thm:projector} supplies a constructive
normal form for the full ideal of convergent Cayley relations and
their shuffle multiples. The exact separating coordinate applies to
that whole ideal. It is stronger than a failure to solve one bounded
matrix system, but still concerns a specified formal presentation.
The conjectural numerical equality for $S_6$ could be proved by
relations that are absent from this ideal. Its $S_8$ analogue is also
untouched. A nonzero formal normal form must not be reported as a
nonzero special value.

\paragraph{Two bookkeeping rules should accompany integration.}
First, rational constant factors in $1-f$ belong to the golden tensor
calculation: the factor $2$ in $1-(1-t)/(1+t)$ is a simple example.
Removing all constant factors would remove logarithmic information.
Second, Cayley-compatible endpoint projection must use the shifted
generators $z+b/2$ and $c-b/2$. Setting the unshifted endpoint
generators to zero does not commute with the involution. These are
requirements of the new proofs, not allegations that the baseline
already made either mistake.

The accompanying \src{integration/INTEGRATION.md} gives concrete
source locations, replacement prose, and the label mapping. All
integration should be compared with the pinned baseline: later
repository changes may already incorporate some of these statements.

\subsection{Research questions on uniform bounds and asymptotics}

\paragraph{1. Optimize the later Euler steps.}
For each integer $N\ge2$, define
\[
 C_N^{\mathrm{av}}=
 \sup_{a,b>0}\E[\mathcal K_N(X_a,Y_b)],\qquad
 C_N^{\mathrm{pt}}=
 \max_{0\le x,y\le1}\mathcal K_N(x,y),
\]
where the boundary value at $y=1$ is continuous. The present theorem
gives $C_N^{\mathrm{av}}\le C_N^{\mathrm{pt}}<9/8$. Determine either
sequence sharply. The two optimization problems should be kept
separate: the gamma laws form a restricted family of probability
measures, so a pointwise maximizer need not be approachable by these
averages. A useful first target is an exact algebraic description of
the pointwise maximum for $N=2$, obtained from its rational stationary
equations and certified boundary comparisons.

\paragraph{2. Determine the fixed-parameter tail asymptotics.}
For fixed $a,b>0$, the exact kernel representation suggests a sharper
description than a universal constant times $2^{-N}$. The transition
of $(1-x^2y^2)^N$ occurs near $xy\asymp N^{-1/2}$. The product
$XY$ has $-\log(XY)$ gamma distributed with shape $a+b$, so Mellin
analysis of that region is natural. The difficulty is to retain the
cancellation in $(W_N(xy)-W_N(x))/(1-y)$ uniformly near $y=1$.
A proof should state its leading term and remainder uniformly on
compact positive parameter sets, then study what fails as an order
approaches zero. This could give much smaller parameter-sensitive
error budgets without weakening the unconditional universal bound.

\paragraph{3. Extend the positive-tail method to other arguments.}
At the Gaussian point, the alternating series and the rational
denominators align particularly well. Which roots of unity, or which
fixed angles in the slit plane, admit a transform whose remainder is
a positive probability average? The first concrete experiment should
be a residue-class decomposition at a primitive sixth root. A useful
theorem would identify the angle-dependent kernel and its sign before
attempting a global numerical constant. At a general angle, positivity
may fail even though the analytic continuation remains well defined.

\subsection{Research questions on radial geometry}

\paragraph{4. Prove or refute global uniqueness of the quartic transition.}
The precise unresolved assertion is that
$Q(a_*(b),b)=0$ has only the certified solution for $0<b<1$.
The determinant calculation already provides a positive derivative
along the threshold in a small neighborhood. An effective program
would combine analytic endpoint expansions at $b\downarrow0$ and
$b\uparrow1$ with a finite interval covering of the remaining compact
interval. A plot of the threshold is not sufficient, especially near
the fully constant point $(1,1)$, where several coefficients vanish.

\paragraph{5. Make the double-turning neighborhood effective.}
The local theorem proves the existence of a fixed radius and a
parameter neighborhood but does not print their sizes. One can seek
explicit bounds for $U_{tt}$ and the fourth and higher Taylor
remainders on a rational parameter box and $0\le t\le T$. The
all-orders recurrence isolates each coefficient, while a Cauchy bound
on the normalized implicit equation can control the infinite tail.
The resulting certificate should exhibit actual rational parameter
pairs with two certified turning-radius intervals. This is a natural
next strengthening of the current existence theorem.

\paragraph{6. Understand the all-orders degeneracy at $(1,1)$.}
At that point $F_{1,1}(z)=\tfrac12\log^2(1-z)$ and the normalized
radius is identically $1/2$. The nearby double-turning transition has
a nonzero sextic coefficient. Determine the first-order parameter
variation of the entire function $\Phi$, rather than one Taylor
coefficient at a time, and ask whether it constrains the number of
small-radius oscillations. This may clarify the relationship between
the isolated certified transition and the infinitely degenerate
constant curve. It must also be reconciled with the still-open
full-radius motion conjectures for integer outer order.

\subsection{Research questions on identities and relation spaces}

\paragraph{7. Lift the golden weight-six and weight-seven candidates.}
The two higher reconstructions in the manuscript are concrete next
targets. Search for a rational-function tensor certificate of the
correct weight whose specialization is exactly the proposed vector,
including its logarithmic completion. Gangl's work
\cite{Gangl} gives relevant higher-weight functional identities and
uses richer multiplicative alphabets. A search restricted to
$t,1-t,1+t$ should therefore be treated as one finite model. The
essential deliverables are an exact tensor row, an endpoint constant,
and ordinary-polylogarithm reconstruction, not only a nullspace rank.
The weight-eight and weight-nine formulas require a further step and
should not be assumed to follow once weights six and seven are solved.

\paragraph{8. Simplify and classify the functional certificates.}
Can the $58$- and $60$-term rows be expressed as short combinations of
classical Kummer specializations? Such a decomposition could reveal
why the golden specialization eliminates the unwanted powers. A
separate finite problem is to determine the exact image of the full
degree-bounded tensor kernel under golden specialization. Its rank
and nullspace can be certified by rational linear algebra. Any
minimality result must explicitly retain the alphabet and degree
bound; it would not establish independence of the resulting periods.

\paragraph{9. Specialize the new functional identities in other fields.}
The identities \eqref{gold:eq:new-functional-identities} hold for all
$0<t<1$, so their use is not confined to the golden ratio. For
example, with $\sigma=\sqrt2-1$, one has
$1-\sigma=\sqrt2\,\sigma$ and $1+\sigma=\sqrt2$. Every table argument
then becomes
\[
 s\,\sigma^{a+b}2^{(b+c)/2}.
\]
Substitution gives two exact weight-five relations in
$\mathbb Q(\sqrt2)$ with the same displayed endpoint constants.
The research problem is to combine them with duplication, inversion,
and other field relations until a small useful basis emerges. This
provides a structured source of new candidate evaluations. It avoids
guessing a large unrelated list of numerical constants at the outset.

\paragraph{10. Add the missing relations to the Cayley quotient.}
For a proposed collection of additional relation vectors
$\mathcal T$, the exact next computation is to ask whether
\[
 \Pi(F)\in\operatorname{span}_{\mathbb Q}\{\Pi(T):T\in\mathcal T\}.
\]
The equivalence with membership in the original Cayley ideal plus
these additional rows follows from the kernel theorem. Thus every
candidate stuffle, distribution, or complementary-argument identity
can be projected immediately. Start with convergent relations and
lower-weight products, then introduce regularized relations only with
their analytic conventions stated explicitly. The nonzero coordinate
$-166348800$ identifies information that the added relations must
actually supply. If the enlarged presentation still fails, report
that precise obstruction; it would not settle numerical independence.

\paragraph{11. Make the quotient computationally practical in higher weight.}
The explicit word projector is finite, but its expanded support grows
quickly. A representation by Cayley-positive Eulerian generators can
retain products in factored form and postpone full shuffle expansion.
Determine the cost of converting a specified target to these
coordinates, and compare it with elimination in a relation matrix.
A useful implementation should return a reconstruction certificate,
so that compression never hides the formal identity it represents.
The all-weight separating-coordinate lemma suggests that many target
tests may need only a few coordinates of the projector.

\paragraph{12. Formalize the smallest complete proof path.}
The Euler result offers a manageable entry point: polynomial
identities, rational Bernstein positivity, the probability-kernel
formula, and the truncated-power mixture. The golden result offers a
different path through finite tensor identities and one-variable
calculus. A proof-assistant project should first choose one complete
analytic theorem, rather than merely formalizing its arithmetic
receipt. In either case, the present records separate the finite
equalities from the mathematical implications they support.

\subsection{Reproduction and the boundary of the deliverable}

The archive is arranged so that the proofs and their finite premises
can be inspected independently. The article contains all definitions,
the golden coefficient tables, the three Euler polynomials, and the
local radial coefficient formulas. The data files contain the full
Bernstein partitions, the rational interval bounds, the exact golden
rows and their specialization receipts, and the Cayley normal form.
The independent review notes explain the delicate analytic and
algebraic steps.

From the extracted directory, the commands
\begin{verbatim}
python3 -m pip install -r requirements.txt
python3 code/run_all.py
latexmk -pdf -interaction=nonstopmode -halt-on-error article.tex
\end{verbatim}
replay the certificates and rebuild the article. The requirements file
lists the Python dependencies used for replay and figure generation;
the two core Euler and radial producers use only Python's standard
library. The golden verifier and independent Euler reviewer use SymPy
for exact symbolic arithmetic. The Cayley verifier uses rational
arithmetic throughout. Figure generation is optional and is invoked
by adding \texttt{--figures} to the replay command.

No integer-relation search is needed to verify the accepted golden
identities, and no numerical root solver is needed to certify the
radial rectangle. The figures visualize the proved structure and are
explicitly labeled as diagnostic or as a limiting local model. The
open questions above remain open in this deliverable; they are not
silently promoted by successful numerical comparisons or by the
strength of a neighboring theorem.
